\begin{ExerciseBlock}
\ExerciseHeading{01}{Statistische und systematische Messabweichungen}{laboratory_measurements_measurement_error_types}
Ordnen Sie die Messabweichungen in den folgenden vier Szenarien ein. Geben Sie jeweils an, wie sich die Abweichung verringern oder korrigieren lässt.
\begin{enumerate}
\item Die Länge eines Metallstabs wird mit einem Zollstock gemessen, der sich dabei durchbiegt.
\item Fünf Personen messen mit einer Stoppuhr die Schwingungsdauer eines gleich ausgelenkten Pendels. Die Ergebnisse streuen.
\item Ein Thermometer zeigt aufgrund einer Fehlkalibrierung konstant zwei Grad Celsius zu viel an.
\item Kleine, wechselnde Luftströmungen beeinflussen die Anzeige einer Balkenwaage.
\end{enumerate}
\end{ExerciseBlock}
\begin{Solution}
\par\medskip\noindent\textbf{Musterlösung}\par\nopagebreak[4]
\begin{enumerate}
\item Systematisch, wenn der Zollstock bei jeder Messung ähnlich durchbiegt. Einen starren Messstab verwenden oder den Zollstock gerade abstützen.
\item Statistisch, soweit die Streuung aus schwankenden Reaktionszeiten stammt. Mehrfach messen und mitteln oder die Zeit automatisch erfassen.
\item Systematisch. Das Thermometer neu kalibrieren oder die bekannte Anzeigeabweichung abziehen.
\item Statistisch bei wechselnden Luftströmungen ohne bevorzugte Richtung. Die Waage abschirmen und Wiederholmessungen durchführen.
\end{enumerate}
\end{Solution}
